%%%%%%%%%%%%%%%%%%%%%%%%%%%%%%%%%%%%%%%%%%%%%%%%%%%%%%%%%%%%%%%%%%%%%%%%%%%%%%%%
%%  GESAMTREFLEXION  (nach dem Hauptteil schreiben)   Umfang: 3–6 Seiten
%%  Gewicht 20 % der Note.
%%%%%%%%%%%%%%%%%%%%%%%%%%%%%%%%%%%%%%%%%%%%%%%%%%%%%%%%%%%%%%%%%%%%%%%%%%%%%%%%
\chapter{Gesamtreflexion}
\label{chap:gesamtreflexion}

\begin{hinweise}
In der Gesamtreflexion \cite{DK2008} überdenken Sie nochmals die Vorlesung und Ihr Portfolio. Sie sollten als
Ausgangspunkt Ihrer Überlegungen Ihr am Anfang ausgefülltes Motivationsformular verwenden (Anhang). Sie sollten
auf folgendes eingehen:
\begin{itemize}\itemsep0pt
\item Welche Erkenntnisse haben Sie aus den Portfolioeinlagen zu Ihrem Projekt gewonnen?
\item Welche Erkenntnisse haben Sie zum Schwerpunktthema gewonnen?
\item Was haben Sie zu Ihrer Art zu Lernen gelernt? Was war förderlich oder hinderlich?
\item Wie hat die Portfolioarbeit Ihre Arbeitsorganisation beeinflusst?
\item Was haben Sie an Allgemeinwissen dazugewonnen?
\item Was haben Sie über sich selber gelernt?
\item Wie beurteilen Sie Ihre Interaktion mit den Kommilitoninnen und Kommilitonen bei der Portfolioerstellung?
\item Wie beurteilen Sie ihren Studien-/Lernerfolg? Wie lautet ihre mit kritischer Bewertung? Welche Schlüsse ziehen Sie für Ihre weitere wissenschaftliche Bildung?
\item Welche Konsequenzen ziehen Sie aus den neu gewonnenen Erkenntnissen?
\item Welche \textbf{neue individuell bedeutsame Fragestellung} (oder Fragestellungen) hat sich aus den neu gewonnenen Erkenntnissen ergeben?
\item Was sind Ihre nächsten Schritte zur Auseinandersetzung mit den neu gewonnenen Fragestellungen?
\item Was ist Ihr Fazit?
\end{itemize}
Bitte schreiben Sie Ihren Text so, dass die geforderten Aussagen identifizierbar sind! Antworten aus nur einem
Satz sind meistens schwer bis unverständlich.
\textbf{Die Gesamtreflexion sollte mindestens drei Seiten und maximal sechs Seiten lang sein. Die Längenangaben
sind ohne Titel, Bilder, Tabellen und Tabellen- oder Bildlegenden zu verstehen.}

\smallskip\emph{Bewertungskriterien zusätzlich beachten:} begründete Einschätzung, inwieweit die
Auseinandersetzung mit dem Schwerpunktthema gelungen ist; fördernde/hemmende Faktoren; was bei Projekt- und
Schwerpunktthema geholfen hat; Ideen zur Fortführung der Forschung und welche Schritte dafür nötig wären;
Konsequenzen für Studium/Beruf. Wird \cite{DK2008} im eigenen Text nicht zitiert, erscheint es auch nicht im
Literaturverzeichnis.
\end{hinweise}

\begin{umfang}{Gesamtreflexion}{3}{6}

\section*{Erkenntnisse zu Projekt- und Schwerpunktthema}
\platzhalter{Rückblick: physikalische Erkenntnisse aus den Einlagen, bezogen auf \Schwerpunkt}

\section*{Lernprozess, Arbeitsorganisation und Zusammenarbeit}
\platzhalter{Art zu lernen, förderlich/hinderlich, Arbeitsorganisation, über sich selbst,
Interaktion mit Kommilitoninnen und Kommilitonen}

\section*{Bewertung des Lernerfolgs}
\platzhalter{kritische Bewertung, Vergleich mit dem Motivationsformular, Schlüsse für die weitere Bildung}

\section*{Neue Fragestellungen und nächste Schritte}
\platzhalter{mindestens eine neue individuell bedeutsame Fragestellung + begründete nächste Schritte}

\section*{Fazit}
\platzhalter{Fazit}

\end{umfang}
